\documentclass[12pt]{article}
\usepackage[utf8]{inputenc}
\usepackage[T1]{fontenc}
\usepackage[margin=1in]{geometry}
\usepackage{setspace}
\usepackage{xcolor}
\usepackage[colorlinks=true, urlcolor=blue!60!black, linkcolor=blue!60!black]{hyperref}
\setlength{\parindent}{0em}
\setlength{\parskip}{0.45em}
\renewcommand{\baselinestretch}{1.0}
\begin{document}
\begin{center}
{\Large \textit{Brown v. Board of Education (Brown II)}, 349 U.S. 294 (1955)\par}
\vspace{0.3em}
{\normalsize Case Brief --- One-Page Edition\par}
\end{center}

\textbf{Memory Jogger.} School boards, supervised by federal equity courts, must admit Black students on a racially nondiscriminatory basis ``with all deliberate speed'' (\href{https://supreme.justia.com/cases/federal/us/349/294/}{Justia}).

\textbf{Facts.} After \textit{Brown I} (1954) held school segregation unconstitutional, the Court heard reargument on relief in five cases from Kansas, South Carolina, Virginia, the District of Columbia, and Delaware. Some districts had begun desegregating; South Carolina and Virginia had done nothing, and Southern leaders were mobilizing massive resistance. NAACP counsel led by Thurgood Marshall urged firm, prompt deadlines; the Attorney General and six Southern states also argued.

\textbf{Procedural History.} Federal courts below denied relief; Delaware's courts ordered immediate admission. The Court reversed the four federal cases and affirmed Delaware.

\textbf{Votes.} Unanimous, 9--0, Chief Justice Warren writing.

\textbf{Holding.} All laws requiring or permitting segregation must yield to \textit{Brown I}. School authorities bear primary responsibility; trial courts must require a ``prompt and reasonable start,'' put the burden of justifying delay on defendants, retain jurisdiction, and order admission ``with all deliberate speed.''

\textbf{Reasoning.} The remaining issue was remedy. Trial courts near local conditions would judge boards' good faith. Equity would balance the plaintiffs' interest in prompt admission against the public interest in an orderly transition, weighing the school plant, transportation, personnel, and attendance areas---but constitutional principles could not yield to disagreement. Only after a prompt start could defendants obtain more time, and only by proving it necessary in the public interest and consistent with good-faith compliance at the earliest practicable date. Warren wrote for a unanimous Court, with no separate opinions.

\textbf{Examples.} \emph{Same-Side:} a prompt one-year plan to redraw bus routes and reassign teachers; a staged prison-reform decree with retained jurisdiction. \emph{Opposite-Side:} delay sought because of public hostility; indefinite tokenism. \emph{Fence-Sitter:} a twelve-year grade-a-year plan that starts now but bypasses older students.

\textbf{Critique.} Entrusting desegregation to segregating boards under a deadline-free formula invited a decade of evasion, until \textit{Alexander v. Holmes County} (1969) repudiated the standard. Originalists add that plaintiffs deserved immediate relief and that the decision launched open-ended structural injunctions.

\textbf{Key Quotation.} ``Traditionally, equity has been characterized by a practical flexibility in shaping its remedies and by a facility for adjusting and reconciling public and private needs.'' --- Chief Justice Warren, 349 U.S. at 300.

\end{document}
